\section*{Bab \ref{ch:g5:geometry} --- Lukisan Geometri}

\begin{solution}{exo:g5:geometry:1}
Kedua lukisannya mengikuti \cref{met:g5:geometry:basic}; \omterm{def:g4:geometry:def}{garis tegak
lurus} dan \omterm{def:g4:geometry:def}{garis sejajar} melalui $A$ itu sendiri saling \omterm{def:g4:geometry:def}{tegak lurus}.
\end{solution}

\begin{solution}{exo:g5:geometry:2}
Rangkaian \omterm{def:g4:geometry:def}{garis tegak lurus} seperti pada \cref{ex:g5:geometry:chain}.
Kalau ketiga sudut yang dibangun tepat siku-siku dan panjangnya
tepat, sudut keempat tidak punya pilihan selain siku-siku juga ---
persegi panjang yang dibangun dengan baik memeriksa dirinya sendiri.
\end{solution}

\begin{solution}{exo:g5:geometry:3}
Alas $[AB]$ sepanjang $7$ cm; busur berjari-jari $6$ cm (dari $A$)
dan $3$ cm (dari $B$); perpotongannya adalah $C$.
\end{solution}

\begin{solution}{exo:g5:geometry:4}
Sama kaki: alas $5$ cm, kedua busurnya dengan bukaan $7$ cm yang
sama. Sama sisi: alas $6$ cm, kedua busurnya dengan bukaan $6$ cm.
\end{solution}

\begin{solution}{exo:g5:geometry:5}
Penggaris siku menjamin \omterm{def:g3:shapes:rightangle}{sudut siku-sikunya}; jangka (atau penggaris)
membawa sisi $5$ cm itu ke keempat sisinya.
\end{solution}

\begin{solution}{exo:g5:geometry:6}
Gambarlah sisi $[AB]$ sepanjang $4$ cm, pilihlah arah untuk $[AD]$
(jangan \omterm{def:g4:geometry:def}{tegak lurus}, agar bukan persegi), dengan $AD = 4$ cm; lalu
busur $4$ cm dari $B$ dan dari $D$ berpotongan di $C$. Jangka
membenarkan keempat sisinya sama panjang.
\end{solution}

\begin{solution}{exo:g5:geometry:7}
Lingkaran berpusat $A$ dengan \omterm{def:g3:shapes:shapes}{jari-jari} $4$ cm dan lingkaran
berpusat $B$ dengan \omterm{def:g3:shapes:shapes}{jari-jari} $3$ cm berpotongan di \emph{dua} titik
(satu di atas, satu di bawah garis $(AB)$), karena $4 + 3 > 5$ dan kedua lingkaran itu bertumpang tindih.
\end{solution}

\begin{solution}{exo:g5:geometry:8}
Satu program yang benar: ``1. Gambarlah persegi panjang $ABCD$
dengan $AB = 6$ cm (alas) dan $BC = 2$ cm. 2. Gambarlah dua busur
berjari-jari $4$ cm yang berpusat di $D$ dan di $C$ (pojok atasnya);
namai $S$ perpotongannya di atas persegi panjang itu. 3. Gambarlah
$[DS]$ dan $[CS]$.'' (Ujung pensilnya adalah segitiga sama kaki $DSC$.)
\end{solution}

\begin{solution}{exo:g5:geometry:9}
Busur berjari-jari $3$ cm dan $4$ cm, yang digambar dari kedua ujung
alas $8$ cm, tidak pernah bertemu: bahkan bila disambung ujung ke
ujung, $3 + 4 = 7$ cm tidak sanggup merentang $8$ cm. Dua sisi
bersama-sama harus selalu lebih panjang daripada sisi ketiganya ---
kalau tidak, segitiganya tidak dapat menutup.
\end{solution}

\begin{solution}{exo:g5:geometry:10}
Gambarlah ruas garis $6$ cm beserta titik tengahnya $O$; gambarlah
\omterm{def:g4:geometry:def}{garis tegak lurus} di $O$; tandai padanya dua titik yang berjarak
$2$ cm dari $O$ (satu di tiap sisi); lalu hubungkan keempat ujungnya
berurutan. Jangka membenarkan keempat sisinya sama panjang.
\end{solution}

\begin{solution}{exo:g5:geometry:11}
Pada denahnya, kambing itu mencapai cakram berpusat $A$ dengan
\omterm{def:g3:shapes:shapes}{jari-jari} $4$ cm, yang terpotong oleh garis pagar pada jarak $3$ cm
dari $A$: daerahnya adalah cakram itu dikurangi irisan di seberang
pagar. Di sepanjang pagar ia dapat merumput pada sebuah ruas garis
(tempat lingkarannya memotong garis pagar); di seberangnya, tidak ada.

\end{solution}
